\section{Experiments}
\label{sec:experiments}

\subsection{Experimental Setup}
\label{sec:setup}

\paragraph{Datasets.}
We evaluate on ImageNet-100 (both 10-class and 100-class subsets)~\cite{russakovsky2014imagenet}, ImageNette, ImageWoof, a 200-class ImageNet subset, and the full ImageNet-1K (1000 classes). To test generalization beyond ImageNet, we evaluate on Food-101~\cite{bossard2014food101} (101 food categories) and CIFAR-10~\cite{krizhevsky2009learning} (10 object categories) using Stable Diffusion~\cite{rombach2021highresolution} 1.5 as the generator. All images are resized to $256 \times 256$ pixels.

\paragraph{Architecture and protocol.}
Following MGD$^3$~\cite{chansantiago2025modeguided}, we use a pretrained DiT-XL/2~\cite{peebles2022scalable} as the frozen generator for all ImageNet-derived experiments, with 50 denoising steps and a high-noise guidance window covering the first 25 timesteps. We evaluate downstream classification using ConvNet-6, ResNet-18, and ResNetAP-10~\cite{kaiming2015deep}, all with instance normalization. Classifiers are trained with SGD (lr $0.1$, weight decay $1\text{e}{-4}$, batch size $128$) with a multi-step LR schedule. We use 2000 epochs for 100-class experiments, ImageNette, and ImageWoof; 3000 for the 10-class subset and CIFAR-10; 1000 for IN-1K. Data augmentation includes random resized crop, horizontal flip, color jitter, PCA lighting noise, and CutMix. Following MGD$^3$, we report the \emph{best} top-1 accuracy across all training epochs. All experiments use three random seeds, and we report mean $\pm$ standard deviation (population std for reporting; sample std for $t$-tests).

\paragraph{Baselines.}
We compare against: (1) \textbf{Unguided} ($\lambda = 0$), the MGD$^3$ baseline; (2) \textbf{Fixed $\lambda$}, globally fixed guidance strength; and (3) \textbf{Random Herding}, randomly selecting IPC real images per class. We additionally compare with D4M~\cite{su2024dataset} under the same DiT generator.

\paragraph{\ags{} configurations.}
For \cags{}, we use the complexity weighting $(\alpha, \beta, \gamma, \delta) = (0, 0.5, 0.5, 0)$---the optimal configuration identified by our factor ablation (\Cref{sec:factor-ablation}) and weight grid search (Appendix~\ref{sec:weight-learning})---with sigmoid parameters $s = 3.0$, $c_0 = 0.6$, and guidance range $[0.0, 0.06]$.

\subsection{Main Results: ImageNet-100 (100-class)}
\label{sec:results-100cls}

\Cref{tab:results-100cls} presents results on the full 100-class subset with ConvNet-6. The unguided baseline achieves $22.79 \pm 0.18\%$. Fixed guidance at $\lambda = 0.10$ yields $21.31 \pm 0.42\%$ ($-6.5\%$ relative), confirming that a globally fixed guidance strength is suboptimal when class complexity varies across 100 classes: a single $\lambda$ over-guides simple classes while under-guiding complex ones. The original four-factor \cags{} achieves $25.29 \pm 0.27\%$ ($+10.9\%$ relative), while the optimized two-factor configuration $(\alpha{=}0, \beta{=}0.5, \gamma{=}0.5, \delta{=}0)$ achieves $25.70 \pm 0.07\%$ ($+12.8\%$ relative), a further $+0.41\%$ improvement with $3.9\times$ lower variance. The extremely tight standard deviation ($\pm 0.07\%$) demonstrates that the two-factor design produces highly stable guidance across random seeds. On the 10-class subset, \cags{} is neutral ($-0.7\%$ relative at IPC=10), confirming that per-class adaptation provides no benefit when class complexity variation is limited (Appendix~\ref{sec:results-10cls}).

\begin{table}[t]
\centering
\small
\caption{ImageNet-100 (100-class) results with ConvNet-6, IPC=10, 3 seeds. The optimized two-factor \cags{} $(0, 0.5, 0.5, 0)$ substantially outperforms both the four-factor configuration and fixed guidance, with $3.9\times$ lower variance.}
\label{tab:results-100cls}
\begin{tabular}{lcccc}
\toprule
\textbf{Method} & \textbf{Avg. $\lambda$} & \textbf{Top-1 (\%)} & \textbf{Top-5 (\%)} & \textbf{$\Delta$ vs.\ unguided} \\
\midrule
Random Herding & --- & $14.06 \pm 0.48$ & $32.31 \pm 1.06$ & $-8.73$ \\
Unguided & 0.000 & $22.79 \pm 0.18$ & $43.67 \pm 0.07$ & --- \\
Fixed $\lambda = 0.10$ & 0.100 & $21.31 \pm 0.42$ & $41.99 \pm 0.33$ & $-1.48$ \\
\midrule
\cags{} (4-factor) & $\sim$0.042 & $25.29 \pm 0.27$ & $47.11 \pm 0.81$ & $+2.50$ \\
\cags{} (2-factor, optimal) & $\sim$0.042 & $\mathbf{25.70 \pm 0.07}$ & $\mathbf{47.72 \pm 0.35}$ & $\mathbf{+2.91}$ \\
\bottomrule
\end{tabular}
\end{table}

\subsection{Cross-Architecture Evaluation}
\label{sec:cross-arch}

To verify that improvements are architecture-agnostic, we evaluate on ResNet-18 and ResNetAP-10 in addition to ConvNet-6 (\Cref{tab:cross-arch}). Both \cags{} configurations improve over unguided across all three architectures. The stronger gains on deeper architectures (ResNet-18: $+15.1$--$17.6\%$, ResNetAP-10: $+10.8$--$15.2\%$) compared to ConvNet-6 ($+12.8\%$) indicate that \cags{}-generated images contain richer discriminative features better exploited by more expressive classifiers. Cross-architecture consistency is further confirmed on ImageNette and ImageWoof across all three architectures, with gains ranging from $+3.5\%$ to $+14.8\%$ relative (Appendix~\ref{sec:results-nette-woof}).

\begin{table}[t]
\centering
\small
\caption{Cross-architecture evaluation on ImageNet-100 (100-class) with IPC=10, 3 seeds, 2000 epochs. Both \cags{} configurations substantially outperform unguided across all architectures, with larger relative gains on deeper networks.}
\label{tab:cross-arch}
\begin{tabular}{lcccc}
\toprule
\textbf{Architecture} & \textbf{Unguided} & \textbf{\cags{} (4-factor)} & \textbf{\cags{} (2-factor)} & \textbf{Best $\Delta$ (rel.)} \\
\midrule
ConvNet-6 & $22.79 \pm 0.18$ & $25.29 \pm 0.27$ & $\mathbf{25.70 \pm 0.07}$ & $+12.8\%$ \\
ResNet-18 (inst.\ norm) & $26.49 \pm 0.42$ & $\mathbf{31.14 \pm 0.34}$ & $30.50 \pm 0.07$ & $+17.6\%$ \\
ResNetAP-10 (inst.\ norm) & $26.65 \pm 0.16$ & $\mathbf{30.70 \pm 0.07}$ & $29.52 \pm 0.08$ & $+15.2\%$ \\
\bottomrule
\end{tabular}
\end{table}

\subsection{Complexity Factor Ablation}
\label{sec:factor-ablation}

To understand the contribution of each complexity factor, we conduct a factor ablation on ImageNet-100 (100-class, IPC=10, ConvNet-6, 3 seeds) testing single-factor variants, equal weights, and variance-dominated weights (\Cref{tab:factor-ablation}). All factor variants improve over unguided ($22.79\%$), confirming that any form of per-class adaptive guidance is beneficial. The ranking reveals two key patterns. First, cluster entropy alone achieves the highest mean accuracy ($25.31 \pm 0.62\%$), but with the highest seed variance. Second, inter-class separability alone is the weakest ($23.75 \pm 0.09\%$), producing near-constant per-class guidance (std $0.006$)---effectively a fixed guidance at $\lambda \approx 0.04$. Crucially, all three configurations that reduce or eliminate separability weight outperform the full four-factor design ($24.24\%$), while the separability-dominated configuration matches it. This provides strong evidence that separability is the least effective factor and that the degree of per-class guidance variation---not the number of factors---is the key driver of \cags{}'s improvement.

\begin{table}[t]
\centering
\small
\caption{Complexity factor ablation on ImageNet-100 (100-class, IPC=10, ConvNet-6, 3 seeds). Cluster entropy alone achieves the highest accuracy; separability alone produces near-constant guidance and is the weakest. The $\lambda$ std column measures the degree of per-class adaptation.}
\label{tab:factor-ablation}
\begin{tabular}{lcccccc}
\toprule
\textbf{Configuration} & $\alpha$ & $\beta$ & $\gamma$ & $\delta$ & \textbf{Top-1 (\%)} & \textbf{$\lambda$ std} \\
\midrule
Unguided & --- & --- & --- & --- & $22.79 \pm 0.18$ & --- \\
\midrule
\textbf{Entropy only} & 0 & 1 & 0 & 0 & $\mathbf{25.31 \pm 0.62}$ & 0.042 \\
Equal weights & 0.25 & 0.25 & 0.25 & 0.25 & $24.69 \pm 0.29$ & 0.043 \\
No separability & 0.25 & 0.25 & 0.50 & 0 & $24.61 \pm 0.44$ & 0.081 \\
Intra-var only & 0 & 0 & 1 & 0 & $24.46 \pm 0.16$ & 0.154 \\
Sep-dominated & 0.10 & 0.10 & 0.20 & 0.60 & $24.25 \pm 0.44$ & 0.034 \\
Full \cags{} & 0.30 & 0.30 & 0.20 & 0.20 & $24.24 \pm 0.26$ & 0.035 \\
Mode count only & 1 & 0 & 0 & 0 & $23.92 \pm 0.12$ & 0.011 \\
Separability only & 0 & 0 & 0 & 1 & $23.75 \pm 0.09$ & 0.006 \\
\bottomrule
\end{tabular}
\end{table}

The two-factor optimal $(0, 0.5, 0.5, 0)$ resolves this by concentrating all weight on entropy and intra-class variance---the two factors with the widest per-class spread---maximizing the range of guidance strengths. This widens the per-class $\lambda$ variation by $2.3\times$ (std $0.0049$ vs.\ $0.0021$ for the four-factor design), directly translating to larger accuracy gains. We provide detailed analysis of per-class mechanisms, quartile analysis, and the separability paradox in Appendix~\ref{sec:class-analysis}.

\subsection{Statistical Significance}
\label{sec:significance}

Despite using only three seeds, the improvements on the 100-class subset are statistically significant ($p < 0.05$) at all IPC settings, with large effect sizes (Cohen's $d \geq 3.90$). Cross-architecture results are also significant ($p = 0.011$ for ResNet-18, $p < 0.001$ for ResNetAP-10). The neutral 10-class IPC=10 result is confirmed non-significant ($p = 0.721$). The ImageNette improvement reaches significance ($p = 0.014$, $d = 4.83$).

\begin{table}[t]
\centering
\small
\caption{Statistical significance of \cags{} vs.\ unguided (and \iast{}+\cags{} vs.\ \cags{}), using paired two-tailed $t$-tests on 3 seeds. Despite limited statistical power, key results reach significance with large effect sizes.}
\label{tab:significance}
\begin{tabular}{llcccc}
\toprule
\textbf{Setting} & \textbf{Comparison} & \textbf{$\Delta$ (\%)} & \textbf{$t$-stat} & \textbf{$p$-value} & \textbf{Cohen's $d$} \\
\midrule
100-cls IPC=1 & CAGS vs Ung & $+0.72$ & $6.80$ & $0.021$ & $3.93$ \\
100-cls IPC=10 & CAGS vs Ung & $+2.50$ & $15.46$ & $0.004$ & $8.93$ \\
100-cls IPC=50 & CAGS vs Ung & $+5.14$ & $14.13$ & $0.005$ & $8.16$ \\
100-cls IPC=1 & IAST+CAGS vs CAGS & $-0.70$ & $-4.17$ & $0.053$ & $-2.41$ \\
\midrule
ResNet-18 IPC=10 & CAGS vs Ung & $+4.65$ & $9.60$ & $0.011$ & $5.55$ \\
ResNetAP-10 IPC=10 & CAGS vs Ung & $+4.05$ & $48.52$ & $<0.001$ & $28.02$ \\
\midrule
10-cls IPC=1 & CAGS vs Ung & $-2.93$ & $-4.21$ & $0.052$ & $-2.43$ \\
10-cls IPC=10 & CAGS vs Ung & $-0.33$ & $-0.41$ & $0.721$ & $-0.24$ \\
10-cls IPC=50 & CAGS vs Ung & $+5.13$ & $3.42$ & $0.076$ & $1.97$ \\
\midrule
ImageNette & CAGS vs Ung & $+2.65$ & $8.37$ & $0.014$ & $4.83$ \\
\bottomrule
\end{tabular}
\end{table}

\subsection{Comparison with D4M}
\label{sec:d4m-summary}

We implement D4M~\cite{su2024dataset} under our exact protocol (same DiT-XL/2 generator, 2000 epochs, 3 seeds, ConvNet-6). D4M encodes real images through the VAE, adds partial noise at timestep $t_{\text{start}}$, and denoises with the frozen DiT. At $t_{\text{start}}{=}25$, D4M underperforms unguided ($20.21\%$ vs.\ $22.79\%$); at $t_{\text{start}}{=}40$, it falls short of \cags{} ($23.07\%$ vs.\ $25.29\%$). This confirms that \cags{}'s training-free, hyperparameter-free approach is more principled than real-image initialization, which requires tuning $t_{\text{start}}$ and access to real data. Full comparison with published DD methods is in Appendix~\ref{sec:dd-baselines}.
